\section{Problem and Executable Joint Recovery}
\label{sec:problem}

An owner-partitioned conflict graph is $G=(V,\mathcal E,w,\alpha)$,
where $w_v\geq0$ and $\alpha:V\to\{1,\ldots,m\}$ assigns vertices to
agents. Writing $w(U)=\sum_{v\in U}w_v$, a feasible schedule is an
independent set; maximum-weight independent set (MWIS) maximizes $w(S)$
over such sets.
In satellite--ground scheduling, vertices represent contacts, satellite
agents own their contacts, and shared satellite or ground resources induce
conflicts. Ownership adds no one-contact-per-agent constraint or communication
protocol. The public-graph partition $\alpha(v)=1+(v\bmod8)$ is artificial;
it does not encode satellite or ground-station ownership.

Given an independent incumbent $S$, an action $a=(C_a,E_a)$ commits an
independent $C_a\subseteq V\setminus S$ and optionally releases
$E_a\subseteq S$. With $N(U)=\bigcup_{v\in U}N(v)$, define
\begin{align}
 D_a&=(N(C_a)\cap S)\cup E_a,\qquad
 B_a=(S\setminus D_a)\cup C_a,\nonumber\\
 q_a&=w(C_a)-w(D_a).\label{eq:v4-action-base}
\end{align}
The action pool uses eleven target cells:
$|C_a|\in\{0,2,4,8\}$ and release targets $\{0,16,64\}$, excluding
the empty action. Actual releases can be fewer; unavailable and duplicate
actions do not create additional queries. Admissible adjustments span at
least two owner partitions. Pure release actions have $C_a=\emptyset$.

Recovery candidates must be compatible with the entire fixed base:
\begin{equation}
 U_a=(D_a\cup N(D_a))\setminus(B_a\cup N(B_a)),\quad
 R_{a,r}=\operatorname{prefix}_r(U_a).
 \label{eq:v4-recovery-scope}
\end{equation}
The ordering places restorable incumbent vertices first, then other vertices
by decreasing reward and identifier. We fix $r=256$ for both the learned
method and its matched controls. Changing $E_a$ changes the action and base,
rather than the amount of information supplied to one policy. Equal actual
prefixes share one request identity.

\begin{proposition}[Feasible recovery and exact gain]
\label{prop:v4-recovery-decomposition}
For any independent $T\subseteq R_{a,r}$, $B_a\cup T$ is feasible and
\begin{equation}
 w(B_a\cup T)-w(S)=q_a+w(T).
 \label{eq:v4-realized-decomposition}
\end{equation}
In particular, $S\cap R_{a,r}$ is a feasible known recovery.
\end{proposition}
\begin{proof}
Removing $N(C_a)\cap S$ makes $B_a$ independent. Every vertex in $R_{a,r}$
is outside and nonadjacent to $B_a$, so independent $T$ completes it
feasibly. Disjointness gives the weight identity. The incumbent intersection
inherits independence from $S$.
\end{proof}

A request specifies $(a,256,\kappa)$, with finite executor workpoint $\kappa$
in its declared native unit. Every policy first executes the same greedy
portfolio on the recovery graph and compares its actual memberships with
$S\cap R_{a,256}$. Its best feasible membership initializes the warm start
$W$. Earlier, timely same-action recoveries can improve $W$. The executor returns an actual recovery $T$
or a failure; full membership validation and rescoring determine its signed
gain in Eq.~\eqref{eq:v4-realized-decomposition}. Neither finite execution
nor a fractional occupancy estimate is an exact recovery oracle.

For total decision allowance $D$, a policy pays for proposals, scope and
priority preparation, inference, execution, and final validation within the
same decision. Let $S_\pi^{\rm ret}(D)$ be its feasible, nondecreasing
schedule actually received by the caller strictly before $t_0+D$, and set
$S_\pi^{\rm ret}(D)=S$ otherwise. The primary objective is
$g_\pi(D)=w(S_\pi^{\rm ret}(D))-w(S)$.
Internal availability before $D$ is a separate quantity. The target is
executable improvement under a complete return deadline, rather than a sum
of individual replacements or the best recovery found after the deadline.
